\section{Runtime Typing Consequences in Singular Operations}

At this point, the basic pieces are in place: names are bindings, singular values are objects, and operators are dispatched according to object behavior. This section brings those ideas together and states their direct consequence for Python's type system.

Python is dynamic at the level of names. A name can first refer to an integer, later to a string, later to \texttt{None}, and Python will allow that rebinding. But Python is not loose at the level of operations. Once an expression executes, the actual objects involved must support that operation.

This is why the same name can legally move from \texttt{42} to \texttt{"42"}, but the expression \texttt{"42" + 1} still fails. Python does not silently guess whether the programmer wanted numeric addition or string concatenation. The programmer must make the conversion explicit.

The purpose of this section is to make that boundary precise:

\begin{quote}
Names are flexible. Objects are typed. Operations are checked.
\end{quote}

This distinction explains much of Python's everyday behavior: why reassignment is easy, why \texttt{type()} and \texttt{isinstance()} inspect runtime objects, why mixed numerical operations can promote toward \texttt{float} or \texttt{complex}, and why unrelated combinations such as text plus number raise errors instead of being silently coerced.